% !TEX program = xelatex
\documentclass{英语试卷}
\begin{document}
\examtitle{衢州、丽水、湖州 2025 年 4 月三地市高三教学质量检测}{英语}
\par 衢州、丽水、湖州 2025 年 4 月三地市高三教学质量检测\par
\par 英语学科 试题卷\par
\begin{notice}
\par 1. 答卷前，考生务必将自己的姓名、准考证号填写在答题卡上。\par
\par 2. 回答选择题时，选出每小题答案后，用铅笔把答题卡上对应题目的答案标号涂黑。如需改动，用橡皮擦干净后，再选涂其他答案标号。回答非选择题时，将答案写在答题卡上，写在本试卷上无效，\par
\par 3. 考试结束后，将本试卷和答题卡一并交回。\par
\end{notice}
\parttitle{第一部分 听力（共两节，满分 30 分）}
\par 做题时，先将答案标在试卷上，录音内容结束后，你将有两分钟的时间将试卷上的答案转涂到答题卡上。\par
\sectitle{第一节（共5小题：每小题1.5分，满分7.5分）}
\par 听下面 5 段对话。每段对话后有一个小题，从题中所给的 A、B、C 三个选项中选出最佳选项。听完每段对话后，你都有 10 秒钟的时间来回答有关小题和阅读下一小题。每段对话仅读一遍，\par
\question{1. How many tests does the man have to take today?}
\begin{choices}(3)
  \item One.
  \item Two.
  \item Three.
\end{choices}
\question{2. Why is Chloe wearing the boots?}
\begin{choices}(3)
  \item She was hiking.
  \item She follows fashion.
  \item She got them from her friend.
\end{choices}
\question{3. Where are the speakers probably?}
\begin{choices}(3)
  \item In an office.
  \item At home.
  \item On a train.
\end{choices}
\question{4. How does the woman find the movie plot?}
\begin{choices}(3)
  \item Creative.
  \item Unexpected.
  \item Old-fashioned.
\end{choices}
\question{5. What is the relationship between the speakers?}
\begin{choices}(3)
  \item Aunt and nephew.
  \item Uncle and niece.
  \item Cousins.
\end{choices}
\sectitle{第二节（共15小题：每小题1.5分，满分22.5分）}
\par 听下面 5 段对话或独白。每段对话或独白后有几个小题，从题中所给的 A、B、C 三个选项中选出最佳选项。听每段对话或独白前，你将有时间阅读各个小题，每小题 5 秒钟；听完后，各小题将给出 5 秒钟的作答时间。每段对话或独白读两遍。 听第 6 段材料， 回答第 6、7 题。\par
\question{6. What is going on at the school tonight?}
\begin{choices}(3)
  \item A swimming course.
  \item An orchestral concert.
  \item A talent competition show.
\end{choices}
\question{7. What is the man doing?}
\begin{choices}(3)
  \item Recommending a hobby to the woman.
  \item Inviting the woman to see a performance.
  \item Persuading the woman to enter a competition.
\end{choices}
\par 听第 7 段材料，回答第 8 至 10 题。\par
\question{8. What did the speakers experience?}
\begin{choices}(3)
  \item A storm.
  \item An earthquake.
  \item A flood.
\end{choices}
\question{9. Why can't the speakers watch the news?}
\begin{choices}(3)
  \item The TV is damaged.
  \item There is a power failure.
  \item The TV reception is poor.
\end{choices}
\question{10. What's the woman's suggestion for dealing with emergencies?}
\begin{choices}(3)
  \item Get some batteries.
  \item Buy some food supplies.
  \item Make their house stronger.
\end{choices}
\par 听第 8 段材料，回答第 11 至 13 题，\par
\question{11. What is the man worried about?}
\begin{choices}(3)
  \item Designing the rooms.
  \item Living near noisy streets.
  \item Spending too much money.
\end{choices}
\question{12. Where will the speakers look for their home location?}
\begin{choices}(3)
  \item In the countryside.
  \item In the downtown area
  \item By the sea.
\end{choices}
\question{13. What does the man suggest doing?}
\begin{choices}(3)
  \item Consulting friends.
  \item Doing a search online.
  \item Visiting houses in person.
\end{choices}
\par 听第 9 段材料，回答第 14 至 17 题，\par
\question{14. When does the conversation probably take place?}
\begin{choices}(3)
  \item In June.
  \item In August.
  \item In September.
\end{choices}
\question{15. Why did the speakers go to the museum?}
\begin{choices}(3)
  \item They were advised to go there.
  \item They had a history class there.
  \item They followed a tour group there.
\end{choices}
\question{16. What did the woman find the most impressive about the museum?}
\begin{choices}(3)
  \item Its number of pieces.
  \item Its unique design
  \item Its collection of calligraphy.
\end{choices}
\question{17. How will the speakers probably go home?}
\begin{choices}(3)
  \item By bike.
  \item By subway.
  \item On foot.
\end{choices}
\par 听第 10 段材料，回答第 18 至 20 题。\par
\question{18. Why does the speaker give the talk?}
\begin{choices}(3)
  \item To train the team how to run.
  \item To arrange a game for the boss.
  \item To sum up experience from failure.
\end{choices}
\question{19. How many goals did the team score in the last game?}
\begin{choices}(3)
  \item 1.
  \item 5.
  \item 20.
\end{choices}
\question{20. What's the team's main task today?}
\begin{choices}(3)
  \item Doing lots of running.
  \item Learning more football ski ls.
  \item Practising shooting one by one.
\end{choices}
\parttitle{第二部分 阅读理解（共两节，满分 50 分）}
\sectitle{第一节（共15小题：每小题2.5分，满分37.5分）}
\par 阅读下列短文，从每题所给的 A、B、C 和 D 四个选项中，选出最佳答案。\par
\passagelabel{A}
\par Mount Seymour Provincial Park features 3, 508 hectares and offers viewpoints overlooking Burrard Inlet, the city of Vancouver. The park has an extensive trail network, and many of the trails lead to notable park features and spectacular lookouts. Here are things to do and see while you are visiting.\par
\par Hiking\par
\par Mt Seymour has 14 hiking trails of varying difficulty and length, offering excellent wildlife viewing opportunities. Obey posted signs and stay on designated marked trails to decrease erosion and trail deterioration, so walking around freely is not allowed. Boardwalks may be slippery.\par
\par Mountain Biking\par
\par Mountain biking is permitted only on designated trails within the park boundary. Trail maps are available online and on kiosks (亭子) at each day-use area. Riders can enjoy the natural scenery along the trail, and it is the responsibility of the riders to know whether the trail they are riding is legal.\par
\par Fishing\par
\par Elsay Lake, accessed by trail only, can be fished but is not stocked. Sitting by the lake, the beautiful scenery is right in front. And anyone fishing or angling must have an appropriate license and anyone fishing overnight is encouraged to post emergency contact information but fires are not permitted anytime within the park. Winter Recreation\par
\par Winter recreation opportunities include ski touring, downhill skiing and snowboarding. Park visitors interested in sightseeing or recreating in the Controlled Recreation Area managed by Mt Seymour must purchase a pass.\par
\par Mountain weather can change quickly. If you become lost or disoriented, stay where you are until the weather clears or you are found. Never leave the trail! Leave a trip plan with a friend in advance. Enjoy your time at Mount Seymour Provincial Park and help us preserve this beautiful natural space for future generations.\par
\question{21. What can visitors do in Mount Seymour Provincial Park?}
\begin{choices}(2)
  \item Free wandering.
  \item Outdoor barbecue.
  \item Unlicensed fishing.
  \item Downhill skiing
\end{choices}
\question{22. Who is advised to leave emergency contact information?}
\begin{choices}(2)
  \item Anyone who fishes by the Elsay Lake.
  \item Anyone who has an appropriate permit.
  \item Anyone who fishes for the whole night.
  \item Anyone who makes a fire and cooks there.
\end{choices}
\question{23. What do the four activities have in common?}
\begin{choices}(2)
  \item They offer views of nature to visitors
  \item They require all visitors to buy a pass.
  \item They are available throughout the park.
  \item They need to be completed with a friend.
\end{choices}
\passagelabel{B}
\par Keira Knightley, a British actress, has undeniably left an indelible mark on the film industry. With her strikingly delicate yet expressive features, she is famed for roles in films like Pirates of the Caribbean series and Atonement, starting her acting career at a young age and achieving significant success in the film industry. But she may have gotten her start in Hollywood too early.\par
\par At age 16, the London-born actress appeared in the unexpected 2002 box-office hit Bend It Like Beckham. She soon found herself in near-constant demand, in both the U. K. and the U. S. "I had a five-year period between the ages of 17 and 21-ish, and I'm never going to have that kind of success again," says Knightley, now 39. Knightley admits that this period of intense success is unlikely to happen again. She also expresses her sadness about the loss of privacy during her teenage years, "It's very cruel to have your privacy taken away in your teenage years."\par
\par Reflecting on her role in 2003's Love Actually, she considers it inappropriate for her age. In one scene, her newly married character opens the front door to her husband's best friend, played by Andrew Lincoln, then 30, who conveys his love to her with a series of cue cards. Moved as she is, in the end, she rejects his love. While filming the scene, director Richard Curtis told her to stop reacting as if the gesture was inappropriate. "And I'm like, 'but it is quite strange.' Also, I knew I was 17. It only seems like a few years ago that everybody else realized I was 17."\par
\par Having started acting professionally at 6, Knightley is hesitant to let her daughters, aged 9 and 5, follow the same path. Interestingly, they haven't shown an inclination either. "One wants to be a sweet-shop owner," she says. "The other one a bear. I don't think that's going to work, but personally speaking, premature fame is not necessarily a good thing."\par
\question{24. What film made Knightley unexpectedly popular?}
\begin{choices}(2)
  \item Pirates of the Caribbean.
  \item Atonement.
  \item Bend It Like Beckham.
  \item Love Actually.
\end{choices}
\question{25. Why did Knightley feel uncomfortable in her scene in Love Actually?}
\begin{choices}(2)
  \item It misrepresented young love.
  \item It didn't match her acting style.
  \item It stressed the age gap with the actor.
  \item It was a bit mature for her young age.
\end{choices}
\question{26. What made Knightley unwilling for her daughters to be actresses?}
\begin{choices}(2)
  \item The same repetitive life.
  \item Her first-hand experience.
  \item Their lack of acting skills.
  \item The competitive film industry.
\end{choices}
\question{27. Where is the text probably from?}
\begin{choices}(2)
  \item A book review.
  \item A magazine.
  \item A drama script.
  \item An autobiography.
\end{choices}
\passagelabel{C}
\par As the proverb says, a straw shows which way the wind blows. A new study published by the U. S. Centers for Disease Control and Prevention (CDC) shows that the H5N1 bird flu virus is probably circulating undetected in domestic animals in many parts of the country and may be infecting unaware vets (兽医).\par
\par In the health agency's Morbidity and Mortality Weekly Report, a group of researchers from the CDC, the Ohio Department of Health and the American Association of Bovine Practitioners reported the results of analysis they conducted on 150 cow vets from 46 states and Canada. They found that three of them had antibodies for the HSN1 bird flu virus in their blood. However, none of the infected vets recalled having any symptoms, including conjunctivitis, or pink eye, the most commonly reported symptom in human cases.\par
\par Seema Lakdawala, a microbiologist at Emory University in Atlanta, who was not involved in the research, said she was surprised that only 2\% of the vets surveyed tested positive for the antibodies, considering another CDC study showed that 17\% of dairy workers sampled had been infected. But she said she was even more surprised that none of them had known they were infected or that they had worked with infected animals.\par
\par "These surprising results indicate that serum surveillance (血清检测) studies are important to inform risk of infections that are going unnoticed," she said. "Vets are on the front line of the outbreak, and increased biosafety practices like respiratory and eye protection should reduce their exposure risk."\par
\par Jennifer Nuzzo, director of the Pandemic Center at Brown University, described the study as "good and bad news study." "On the one hand, we see concerning evidence that there may be more H5NI outbreaks on farms than being reported," she said, "On the other hand, I'm reassured that there isn't evidence that infections among vets have been widespread."\par
\par However, John Korslund, a U. S. Department of Agriculture scientist, said that finding H5N1 antibodies in the blood of vets was an "interesting but very imprecise way to measure state cattle incidence, so we need follow-up study."\par
\question{28. Why is the proverb in paragraph 1 mentioned?}
\begin{choices}(2)
  \item To describe a fact.
  \item To introduce the study.
  \item To explain an is ue.
  \item To present an opinion.
\end{choices}
\question{29. What is paragraph 2 mainly about?}
\begin{choices}(2)
  \item The findings of the analysis.
  \item The symptoms of the vets.
  \item The causes of H5N1 bird flu.
  \item The basic logic of the study.
\end{choices}
\question{30. What surprised Seema Lakdawala the most?}
\begin{choices}(2)
  \item The infected vets' unawareness of the issue.
  \item The 17\% infection rate among dairy workers.
  \item The high number of positive test results in vets.
  \item The risk of unrecognized infections in cows and vets.
\end{choices}
\question{31. What aspect should be improved in the follow-up study?}
\begin{choices}(2)
  \item The size of the tested group.
  \item The individual samples
  \item The data analysis procedure.
  \item The measuring method.
\end{choices}
\passagelabel{D}
\par OpenAI, the company behind the artificial intelligence tool ChatGPT, has set out restructuring plans to make it easier for it to raise money and remove restrictions made by its non-profit parent (母公司).\par
\par OpenAI reveals its for-profit unit would become a public benefit corporation, a company required to consider the interests of society as well as shareholder (股东) value. The non-profit parent would have "a significant interest" in the public benefit corporation, OpenAI said in a blog-post, adding that it would be one of the "best resourced non-profits in history."\par
\par OpenAI has been looking to make changes to its rules to make it easier for investors to profit when they inject new investment into the firm. OpenAI was founded in 2015 as a research-focused non-profit but created a for-profit unit four years later to secure funding for the high costs of AI development. Its unusual structure gave control of the for-profit unit to the non-profit.\par
\par A S6.6 billion fundraising by OpenAI, which valued it at \$157 billion, was conditional on it being able to remove a profit limit and move towards being a more conventional company.\par
\par "We again need to raise more money than we’d imagined. Investors want to back us but they need conventional equity (股东权益) and a less unique structure," OpenAI said. "The hundreds of billions of dollars that companies are investing into AI show what it will take for OpenAI to continue pursuing the mission."\par
\par "This is the critical step the company needs in order to continue fundraising," said Gil Luria, analyst with DA Davidson \& Co. The key to the announcement, he said, was that the for-profit side of OpenAI "will run and control OpenAI's operations and business." For its AI competitors, it's a different matter. Musk, an OpenAI co-founder, is trying to preclude the plan, as is Meta Platforms, which is developing a competitor against ChatGPT named Meta AI.\par
\question{32. What is a public benefit corporation required to do?}
\begin{choices}(2)
  \item Operate its non-profit parent independently.
  \item Focus only on maximizing shareholder value.
  \item Reduce the amount of money raised from investors.
  \item Combine the interests of society and shareholder value.
\end{choices}
\question{33. Why does OpenAI believe it needs to restructure its operations in paragraph 5?}
\begin{choices}(2)
  \item To attract more investors and simplify the structure.
  \item To push forward its development goals on AI safety.
  \item To settle conflicts with its non-profit parent organization.
  \item To prevent competitors from taking control of its operations.
\end{choices}
\question{34. What does the underlined word "preclude" in paragraph 6 mean?}
\begin{choices}(2)
  \item Design.
  \item Prevent.
  \item Back.
  \item Follow.
\end{choices}
\question{35. Which of the following is the best title for the text?}
\begin{choices}(2)
  \item OpenAI to Keep Influencing Shareholders
  \item OpenAI to Defend ChatGPT from Criticism
  \item OpenAI to Relax Control over Its For-Profit Unit
  \item OpenAI to Spend Funds on Artificial Intelligence
\end{choices}
\sectitle{第二节（共5小题；每小题2.5分，满分12.5分）}
\par 阅读下面短文，从短文后的选项中选出能填入空白处的最佳选项。选项中有两项为多余选项，\par
\par Ugh! You swear you had your keys a second ago, and now they're just gone! Before you panic, take a deep breath. \blanknum{36}. Losing things is a part of life, but finding them doesn't have to be a nerve-wracking rush. In this article, we'll teach you where you can look and how you can retrace your steps to find lost objects.\par
\par Identify what makes your missing object stand out in its environment. It might be its size, color, texture or shape. \blanknum{37} It will make the process faster and more efficient. Instead of looking at everything in a space, this approach helps you focus your attention only on objects with that distinction.\par
\par Look in places where you've lost this item before. Do you tend to lose your keys or student cards often? If so, it may have ended up in a similar place as before. Think about where your common lost objects tend to be lost. For example, maybe you have a habit of leaving your keys in the lock or your bag in the car. \blanknum{38}\par
\par Try looking at the room from a new angle. If you've retraced your steps, gone through all the motions, and still haven't found anything, change your perspective. Crawl instead of walking down the hall. \blanknum{39} Seeing things from a different point of view can help you notice details you may have missed before.\par
\par \blanknum{40} Scattered and crowded areas make it easier to lose objects, as things can fall under other things or get buried beneath piles. Regularly clean and tidy up living spaces, so you won't have to spend more time looking for things later on.\par
\begin{choices}(1)
  \item Or sit where you were standing.
  \item Then search based on that unique feature
  \item Look in the messiest areas of your living space.
  \item Everyone misplaces things every now and again.
  \item Check these places before extending your search.
  \item The best thing you can do to avoid losing things is be organized
  \item Thinking about the last time you saw it helps you retrace your steps.
\end{choices}
\parttitle{第三部分 语言知识运用（共两节，满分 30 分）}
\sectitle{第一节（共15小题；每小题1分，满分15分）}
\par 阅读下面短文，从短文后各题所给的 A、B、C 和 D 四个选项中，选出可以填入空白处的最佳选项，\par
\par George White, 77, may look like an ordinary grandfather at a family reunion—until you learn he calls 200 people his "children". "They're not biologically my kids, but \blanknum{41} they surely are," White told the local news outlet.\par
\par That is how \blanknum{42} he became and still is to the students who \blanknum{43} his school bus, a job he held for 27 years before he retired. Yet, driving a bus was never his \blanknum{44} plan. In 1990, he fell off a roof while working outdoors. After that, he intended to \blanknum{45} a job closer to the ground. However, he said no job has ever \blanknum{46} him higher. "It's the children." White said. "And being in a position where you can love kids every single day is a \blanknum{47} position to be in."\par
\par The positive feeling was \blanknum{48} by so many of the kids on his bus over the years. One of his former students, Kaitlyn Hare, cited that he just made everybody feel safe, loved and cared for. It is a \blanknum{49} so strong that even though George retired years ago, former students \blanknum{50} recently for one last ride. "They're finding their \blanknum{51} that they had 20 years ago." White said. "And now even their children are sitting on their laps. And that kind of feeling is a wonderful thing."\par
\par What was White's secret to developing this \blanknum{52} ? "He only had two \blanknum{53} on the bus," another former student, Louis Castello, said. "Show everyone love, and show respect." It's a lesson many of them now \blanknum{54} with them.\par
\par "I'm convinced that when you love and respect people, that's what you're going to get back." White said. That is how one man's unexpected path became 200 lives' 55.\par
\cz{41}{personally}{emotionally}{mentally}{socially}
\cz{42}{exposed}{committed}{attached}{accustomed}
\cz{43}{drove}{missed}{repaired}{rode}
\cz{44}{realistic}{original}{subsequent}{previous}
\cz{45}{offer}{perform}{seek}{quit}
\cz{46}{held}{pushed}{promoted}{lifted}
\cz{47}{lovely}{tricky}{suitable}{challenging}
\cz{48}{returned}{inspired}{judged}{hidden}
\cz{49}{symbol}{character}{bond}{tone}
\cz{50}{graduated}{gathered}{communicated}{developed}
\cz{51}{seats}{wallets}{puzzles}{jobs}
\cz{52}{strategy}{skill}{affection}{curiosity}
\cz{53}{bans}{options}{records}{rules}
\cz{54}{turn to}{carry through}{keep up}{search for}
\cz{55}{career}{dream}{position}{compass}
\sectitle{第二节（共10小题；每小题1.5ladnepen分，满分15分）}
\par 阅读下面短文，在空白处填入 1 个适当的单词或括号内单词的正确形式。\par
\par On January 28, 3,000 drones (无人机) performed a stunning light show in Chongqing during the Spring Festival Gala. The show, organized by Damoda Intelligent Control Technology, \blanknum{} (feature) patterns of famous landmarks and flowers across the night sky.\par
\par Despite challenges such as unpredictable weather and signal interference, the team, \blanknum{} average age is 27, delivered a flawless performance. Most of the team members specialize in drone technology, while others are enthusiasts recruited by the company.\par
\par In recent years, drone performances \blanknum{58} (become) a major trend in entertainment. Compared to traditional fireworks, drones are \blanknum{59} (safe) and less polluting. Advancements in technology have made larger-scale performances possible. Damoda's third-generation drones are designed \blanknum{60} (take) off from a compact box, saving space, while the fourth-generation models charge automatically. In 2023, the company set two Guinness World Records by \blanknum{61} (fly) 10,197 drones simultaneously.\par
\par Competition is shifting from technology to creativity, as the industry grows, with a greater emphasis \blanknum{62} storytelling and cultural elements. Beyond entertainment, drones are also fueling \blanknum{63} low-altitude economy, with widespread use in logistics, agriculture, and emergency rescue.\par
\par \blanknum{64} is expected by industry experts, drone shows are going to become even more spectacular, with continuous technological progress. China's drone industry continues to grow globally, setting new standards and reinforcing its \blanknum{65} (influence) position in aerial technology.\par
\parttitle{第四部分 写作（共两节，满分 40 分）}
\sectitle{第一节（满分15分）}
\par 假定你是李华，你校有些学生把老师的照片制作成表情包(meme)并在社交媒体上传播，针对这一现象，请给校英语报 Letters 专栏投稿，内容包括：\par
\par (1) 陈述观点;\par
\par (2) 提出建议。\par
\par 注意：（1）写作词数应为80左右；（2）请按如下格式在答题纸的相应位置作答。\par
\par Dear Editor,\par
\sectitle{第二节（满分25分）}
\par 阅读下面材料，根据其内容和所给段落开头语续写两段，使之构成一篇完整的短文。\par
\par Tim had always struggled with his weight. He loved food, but at school, be often found himself the target of cruel jokes and teasing.\par
\par His classmates would call him names like "Big Tim" and make fun of his size, and though Tim tried to ignore it, it still hurt. Every day, he wished he could be like the athletic kids, running fast and playing sports with case, but that always seemed out of reach.\par
\par One day, a group of his classmates—three boys who were notorious for their pranks (恶作剧)—decided to play a trick on him. They knew Tim wasn't athletic, and they thought it would be funny to sign him up for the upcoming school's long-distance race. They figured he'd panic when he found out and that it would be fun to watch him fail. Without telling Tim, they registered his name for the race.\par
\par On race day, as Tim walked across the track trying to join his classmates in the stands, be saw some of them holding giant posters of his face with slogans like "The Big Runner" and "Go Tim!" He froze. He realized what had happened—the boys had tricked him into racing. His face turned red with embarrassment, but the crowd was already watching, and he felt trapped.\par
\par The starting gun fired, and to everyone's surprise, Tim began to jog. It wasn't fast, and he could barely keep up with the other names, but he kept moving. His legs burned with each step, and his breath came in sharp gasps, but he didn't stop. His classmates' laughter echoed in his cars, but he ignored it.\par
\par As the race went on, something shifted. A few students who weren't part of the prank started cheering for him. "Come on, Tim! You can do it!" they yelled. Their encouragement made Tim feel lighter, and his pace picked up just a little. It wasn't about winning anymore; it was about proving to himself that be could finish.\par
\par 注意：\par
\par (1) 续写词数应为 150 左右;\par
\par (2) 请按如下格式在答题卡的相应位置作答。\par
\writingpara{Paragraph 1: By the time Tim reached the final lap, he was tired out.}
\writinglines{8}
\writingpara{Paragraph 2: The crowd erupted in applause.}
\writinglines{8}
\end{document}
